
\documentclass[11pt,a4paper]{article}
\usepackage[utf8]{inputenc}
\usepackage[T1]{fontenc}
\usepackage{lmodern}
\usepackage[a4paper,margin=1in]{geometry}
\usepackage{amsmath,amssymb}
\usepackage{booktabs,longtable,array}
\usepackage{listings}
\usepackage{xcolor}
\usepackage{hyperref}
\usepackage{enumitem}
\usepackage{setspace}
\usepackage{titlesec}

\hypersetup{colorlinks=true,linkcolor=blue,urlcolor=blue,citecolor=blue}
\setstretch{1.12}
\titleformat{\section}{\large\bfseries}{\thesection.}{0.5em}{}
\titleformat{\subsection}{\normalsize\bfseries}{\thesubsection.}{0.5em}{}

\lstdefinestyle{mystyle}{
  basicstyle=\ttfamily\small,
  breaklines=true,
  columns=fullflexible,
  frame=single,
  keepspaces=true,
  showstringspaces=false,
  keywordstyle=\color{blue},
  commentstyle=\color{teal!70!black},
  stringstyle=\color{red!60!black}
}
\lstset{style=mystyle}

\begin{document}

\begin{center}
{\LARGE \textbf{English LaTeX Translation and Consolidation of the MMHCL-Plus Notes}}\\[0.5em]
{\large Spectral Radius Tracking, Mutual Information Diagnostics, and Q1-Style Ablation Design}
\end{center}

\section{Context}
Based on a careful analysis of the GitHub codebase, including \texttt{TwoStageTrainer (Rev5.2)}, \texttt{LayerwiseEncoder} with \texttt{max\_g\_layers=2}, \texttt{build\_hypergraph\_laplacian} with
\[
\Theta = D_v^{-1/2} H W_e D_e^{-1} H^T D_v^{-1/2},
\]
\texttt{svd\_filter\_incidence}, and the configuration \texttt{MMHCLPlusConfig} with five tasks and \(D = 4096\), the following implementation package is proposed for Spectral Radius Tracking and MINE/KSG Mutual Information Estimation as a DPI sanity check.

This consolidated document is written to improve coherence and cohesion relative to the raw notes. It organizes the material into a theoretical foundation, diagnostic modules, ablation variants, notebook integration, and statistical validation.

\section{Theoretical foundation}
\subsection{Why these two tools are needed}
The Data Processing Inequality (DPI) states that if
\[
X \rightarrow h^{(l)} \rightarrow h^{(l+1)}
\]
forms a Markov chain, and each GNN layer is deterministic, then
\[
I(X; h^{(l)}) \ge I(X; h^{(l+1)}).
\]
In other words, mutual information (MI) should decrease monotonically across layers.

At the same time, the spectral radius \(\lambda_{\max}(\Theta)\) of the diffusion operator \(\Theta\) controls the rate of oversmoothing. When \(\lambda_{\max}(\Theta)\) approaches 1, node embeddings drift toward the stationary distribution, which is precisely the oversmoothing regime. Neighbor-Layer Contrastive Learning (Neighbor-Layer CL) is intended to counteract this effect by preserving a meaningful separation between \(h^{(l)}\) and \(h^{(l+1)}\).

\subsection{Sanity-check criteria}
If the architecture is functioning as intended, the following sanity checks should hold:
\begin{itemize}[leftmargin=1.5em]
\item Mutual information should decrease monotonically across layers, consistent with DPI.
\item The spectral radius should remain below 1 and should decrease when Dirichlet regularization is enabled.
\item Neighbor-Layer CL should slow the rate of MI decay relative to a baseline without NLCL.
\end{itemize}

\section{Diagnostic package structure}
\subsection{New file: \texttt{codes/mmhcl\_plus/diagnostics/\_\_init\_\_.py}}
\begin{lstlisting}[language=Python]
from .spectral_radius import (
    compute_spectral_radius,
    track_spectral_radius_per_epoch,
    SpectralRadiusTracker,
)
from .mutual_information import (
    MINEEstimator,
    ksg_mutual_information,
    LayerwiseMITracker,
)
\end{lstlisting}

\subsection{New file: \texttt{codes/mmhcl\_plus/diagnostics/spectral\_radius.py}}
The purpose of this module is to track \(\lambda_{\max}(\Theta)\), where \(\Theta = I - L\) is the diffusion operator derived from the hypergraph Laplacian built by \texttt{build\_hypergraph\_laplacian()}.

The key architectural insight in Rev5.2 is the following:
\begin{itemize}[leftmargin=1.5em]
\item \(L = I - \Theta\), so the eigenvalues of \(\Theta\) are \(1 - \lambda_i(L)\).
\item \(\lambda_{\max}(\Theta)\) controls the rate of oversmoothing.
\item Dirichlet energy regularization pushes the eigenvalues of \(L\) away from 0, which is equivalent to pulling \(\lambda_{\max}(\Theta)\) away from 1.
\end{itemize}

A typical usage pattern is:
\begin{lstlisting}[language=Python]
tracker = SpectralRadiusTracker(interaction_matrix)
for epoch in range(n_epochs):
    # training step
    tracker.log(epoch)
tracker.plot()
tracker.to_dataframe()
\end{lstlisting}

\subsection{Core functions in \texttt{spectral\_radius.py}}
\begin{itemize}[leftmargin=1.5em]
\item \texttt{compute\_spectral\_radius(laplacian, method="arpack", tol=1e-6)} computes the spectral radius of \(\Theta\) without explicitly forming \(\Theta\) for large sparse graphs.
\item \texttt{compute\_spectral\_gap(laplacian, tol=1e-6)} computes \(\lambda_2(L)\), the second-smallest eigenvalue of \(L\), as a measure of connectivity and smoothing behavior.
\item \texttt{compute\_layer\_spectral\_radius(layer\_embeddings)} estimates an empirical transition map between consecutive layer embeddings and computes the corresponding spectral contraction factor.
\item \texttt{track\_spectral\_radius\_per\_epoch(interaction\_matrix, svd\_top\_k=10)} performs a one-shot comparison of the spectral radius before and after SVD filtering.
\end{itemize}

\subsection{Class: \texttt{SpectralRadiusTracker}}
This tracker is intended for direct integration into the training loop. It stores two kinds of history:
\begin{itemize}[leftmargin=1.5em]
\item Laplacian-based history, including spectral radius and spectral gap per epoch.
\item Layer-based history, including empirical contraction factors between consecutive hidden layers.
\end{itemize}

The class provides methods such as \texttt{log\_laplacian()}, \texttt{log\_layer\_contraction()}, \texttt{update\_laplacian()}, \texttt{to\_dataframe()}, and \texttt{plot()} for analysis and visualization.

\section{Q1-style ablation strategy}
\subsection{Ablation rationale}
According to the revision-52 design notes, the MMHCL-Plus architecture should be evaluated through three complementary ablation groups in order to estimate the marginal contribution of each component:
\begin{itemize}[leftmargin=1.5em]
\item \textbf{A-series}: component-level ablations that switch individual modules on or off.
\item \textbf{B-series}: hyperparameter sensitivity studies, including \(g\)-layers, projector dimension, ramp behavior, and hard-negative delay.
\item \textbf{C-series}: comparisons among alternative loss balancers.
\end{itemize}

\subsection{Proposed ablation matrix}
\begin{longtable}{|>{\raggedright\arraybackslash}p{1.3cm}|>{\raggedright\arraybackslash}p{3.6cm}|>{\raggedright\arraybackslash}p{5.0cm}|>{\raggedright\arraybackslash}p{4.1cm}|}
\hline
\textbf{ID} & \textbf{Variant} & \textbf{Description} & \textbf{Scientific goal}\\
\hline
\endfirsthead
\hline
\textbf{ID} & \textbf{Variant} & \textbf{Description} & \textbf{Scientific goal}\\
\hline
\endhead
A0 & Full (MMHCL-Plus) & All five losses, Hybrid Balancer, warmup ramp, delayed FAISS, VICReg with \(D=4096\), and SVD spectral augmentation & Upstream reference baseline\\
\hline
A1 & w/o Neighbor-Layer CL & Disable \(\mathcal{L}_{u2u}\) and \(\mathcal{L}_{i2i}\), effectively reverting toward the original MMHCL & Measure the contribution of the NLGCL principle\\
\hline
A2 & w/o SVD Spectral Aug & Set \texttt{apply\_svd\_filtering=False} & Quantify the contribution of spectral debiasing against popularity bias\\
\hline
A3 & w/o VICReg Expansion & Reduce the projector dimension to 1024 & Test the ``Dimensionality Paradox'' claim\\
\hline
A4 & w/o CL Warmup Ramp & Activate contrastive learning abruptly at warmup & Measure the effect of CL activation shock\\
\hline
A5 & w/o Delayed FAISS & Use hard negatives from epoch 0 & Validate the motivation for a delay of 50 epochs\\
\hline
A6 & w/o Hutchinson Dirichlet & Disable \(\mathcal{L}_{Dir}\) & Measure the anti-oversmoothing contribution of the Dirichlet term\\
\hline
A7 & w/ Ego-Final Anchor & Re-enable \(\mathcal{L}_{ego\_final}\), yielding 6 tasks & Validate the removal of the ego-final anchor in Rev5.2\\
\hline
A8 & w/o Soft BYOL Cross & Disable \(\mathcal{L}_{align}\) & Measure the contribution of cross-branch alignment\\
\hline
B1--B3 & \(g \in \{1,2,3\}\) & Sweep neighbor-layer depth & Validate the DPI-inspired safe zone \(g \in \{1,2\}\)\\
\hline
C1 & Uncertainty only & Kendall-style uncertainty weighting & Baseline loss balancer\\
\hline
C2 & GradNorm only & Chen-style gradient normalization & Baseline loss balancer\\
\hline
C3 & Fixed weights & Set all task weights to 1.0 & Lower-bound baseline with no balancing\\
\hline
\end{longtable}

\section{New file: \texttt{mmhcl\_plus/ablation/ablation\_config.py}}
This file serves as the central ablation registry of the repository.

\begin{lstlisting}[language=Python]
# mmhcl_plus/ablation/ablation_config.py
"""
Ablation configuration registry for MMHCL-Plus (revision 52).
Each AblationVariant toggles a specific architectural component
to measure its marginal contribution on Recall@K / NDCG@K / Precision@K.
"""
from dataclasses import dataclass, field, asdict
from typing import Dict, List, Optional
import json, copy

@dataclass
class AblationVariant:
    name: str
    enable_neighbor_layer_cl: bool = True
    enable_svd_spectral:     bool = True
    vicreg_projector_dim:    int  = 4096
    vicreg_hidden_dim:       int  = 1024
    enable_cl_warmup_ramp:   bool = True
    ramp_epochs:             int  = 20
    enable_delayed_faiss:    bool = True
    delay_hard_negs_epoch:   int  = 50
    enable_dirichlet:        bool = True
    enable_ego_final_anchor: bool = False
    enable_soft_byol_cross:  bool = True
    g_layers: int = 2
    svd_top_k: int = 10
    balancer_type: str = "hybrid"
    balancer_switch_epoch: int = 40
    warmup_epochs: int = 15
    n_tasks: int = 5
    notes: str = ""

    def to_dict(self): return asdict(self)
    def dump(self, path): json.dump(self.to_dict(), open(path, "w"), indent=2)

REGISTRY: Dict[str, AblationVariant] = {
    "A0_full":         AblationVariant(name="A0_full",
                                       notes="Full MMHCL-Plus (rev52)"),
    "A1_no_nlcl":      AblationVariant(name="A1_no_nlcl",
                                       enable_neighbor_layer_cl=False,
                                       n_tasks=3,
                                       notes="Remove Neighbor-Layer CL (back to MMHCL)"),
    "A2_no_svd":       AblationVariant(name="A2_no_svd",
                                       enable_svd_spectral=False,
                                       notes="Disable SVD spectral filtering"),
    "A3_small_proj":   AblationVariant(name="A3_small_proj",
                                       vicreg_projector_dim=1024,
                                       vicreg_hidden_dim=512,
                                       notes="Shrink VICReg D 4096->1024 (Dim. Paradox)"),
    "A4_no_ramp":      AblationVariant(name="A4_no_ramp",
                                       enable_cl_warmup_ramp=False,
                                       notes="Abrupt CL activation (expect shock)"),
    "A5_no_delay":     AblationVariant(name="A5_no_delay",
                                       enable_delayed_faiss=False,
                                       delay_hard_negs_epoch=0,
                                       notes="FAISS hard negatives from epoch 0"),
    "A6_no_dirichlet": AblationVariant(name="A6_no_dirichlet",
                                       enable_dirichlet=False,
                                       n_tasks=4,
                                       notes="Remove Hutchinson Dirichlet"),
    "A7_ego_final":    AblationVariant(name="A7_ego_final",
                                       enable_ego_final_anchor=True,
                                       n_tasks=6,
                                       notes="Re-enable ego-final anchor (paradox)"),
    "A8_no_cross":     AblationVariant(name="A8_no_cross",
                                       enable_soft_byol_cross=False,
                                       n_tasks=4,
                                       notes="Remove Soft BYOL cross-branch"),
    "B1_g1":           AblationVariant(name="B1_g1", g_layers=1,
                                       notes="Neighbor depth g=1"),
    "B2_g2":           AblationVariant(name="B2_g2", g_layers=2,
                                       notes="Neighbor depth g=2 (default)"),
    "B3_g3":           AblationVariant(name="B3_g3", g_layers=3,
                                       notes="Neighbor depth g=3 (beyond DPI safe zone)"),
    "C1_uncertainty":  AblationVariant(name="C1_uncertainty",
                                       balancer_type="uncertainty"),
    "C2_gradnorm":     AblationVariant(name="C2_gradnorm",
                                       balancer_type="gradnorm"),
    "C3_fixed":        AblationVariant(name="C3_fixed",
                                       balancer_type="fixed"),
}

def get(variant_name: str) -> AblationVariant:
    if variant_name not in REGISTRY:
        raise KeyError(f"Unknown ablation variant '{variant_name}'. "
                       f"Available: {list(REGISTRY)}")
    return copy.deepcopy(REGISTRY[variant_name])
\end{lstlisting}

\section{New file: \texttt{mmhcl\_plus/trainer/hybrid\_balancer.py}}
This module extends the balancer so that it can support the entire C-series of ablations: uncertainty weighting, GradNorm, fixed weights, and the default hybrid switch from uncertainty to GradNorm.

\begin{lstlisting}[language=Python]
class HybridLossBalancer(nn.Module):
    """
    mode in {'hybrid','uncertainty','gradnorm','fixed'}
    - hybrid     : Uncertainty (e<switch) -> GradNorm (e>=switch)  [rev52 default]
    - uncertainty: Kendall et al. 2018 homoscedastic weighting
    - gradnorm   : Chen et al. 2018 adaptive gradient normalization
    - fixed      : w_k = 1.0 for all k
    """
    def __init__(self, num_tasks=5, mode="hybrid", switch_epoch=30, alpha=1.5):
        super().__init__()
        self.mode, self.switch_epoch, self.alpha = mode, switch_epoch, alpha
        self.log_vars   = nn.Parameter(torch.zeros(num_tasks))
        self.gn_weights = nn.Parameter(torch.ones(num_tasks))
        self.register_buffer("L0", torch.zeros(num_tasks))
        self._L0_set = False

    def _uncertainty(self, losses):
        prec = torch.exp(-self.log_vars)
        return (prec * losses + self.log_vars).sum()

    def _fixed(self, losses):
        return losses.sum()

    def _gradnorm(self, losses, shared_params):
        if not self._L0_set:
            self.L0.copy_(losses.detach())
            self._L0_set = True
        w = self.gn_weights
        weighted = (w * losses).sum()
        return weighted

    def forward(self, losses: torch.Tensor, shared_params=None, epoch: int = 0):
        if self.mode == "fixed":
            return self._fixed(losses)
        elif self.mode == "uncertainty":
            return self._uncertainty(losses)
        elif self.mode == "gradnorm":
            return self._gradnorm(losses, shared_params)
        elif self.mode == "hybrid":
            return self._uncertainty(losses) if epoch < self.switch_epoch \
                   else self._gradnorm(losses, shared_params)
        raise ValueError(self.mode)
\end{lstlisting}

\section{Notebook insertion cells}
The notes propose appending new cells to the end of \texttt{Local\_Random\_seeds\_train\_mmhcl\_baby\_colab\_original.ipynb}.

\subsection{Cell [Ablation-1]: import and select a variant}
\begin{lstlisting}[language=Python]
# === ABLATION DRIVER CELL ===
import os, json, wandb, torch, numpy as np, random
from mmhcl_plus.ablation.ablation_config import REGISTRY, get
from mmhcl_plus.trainer.hybrid_balancer import HybridLossBalancer

ABLATION_NAME = os.environ.get("ABLATION", "A0_full")
SEED = 612550091
cfg = get(ABLATION_NAME)
print(f"[Ablation] Running variant = {cfg.name} | notes: {cfg.notes}")
print(json.dumps(cfg.to_dict(), indent=2))

random.seed(SEED); np.random.seed(SEED)
torch.manual_seed(SEED); torch.cuda.manual_seed_all(SEED)
torch.backends.cudnn.deterministic = True
\end{lstlisting}

\subsection{Cell [Ablation-2]: patch the model and training pipeline}

\textbf{Important.} Do \emph{not} patch individual fields manually, otherwise
toggles such as \texttt{enable\_neighbor\_layer\_cl}, \texttt{enable\_dirichlet},
\texttt{enable\_soft\_byol\_cross}, \texttt{enable\_ego\_final\_anchor},
\texttt{balancer\_type}, and \texttt{balancer\_transition\_epoch} will silently
stay at their defaults and variants A1 / A6 / A8 / C1--C3 will not actually
differ from A0. Use the centralized \texttt{apply\_to\_args()} helper, which is
the single source of truth for the field-to-argument mapping:

\begin{lstlisting}[language=Python]
from mmhcl_plus.ablation import apply_to_args

apply_to_args(cfg, args)  # patches all 17 fields in one call

projector = ExpandedProjector(
    input_dim=64,
    hidden_dim=args.projector_hidden_dim,
    output_dim=args.projector_out_dim,
).cuda()

hybrid_balancer = HybridLossBalancer(
    num_tasks=args.num_tasks,
    mode=args.balancer_type,
    switch_epoch=args.balancer_transition_epoch,
).cuda()
\end{lstlisting}

When driving the sweep via \texttt{--ablation\_variant} (e.g.\ through the
notebook's Cell~20 or \texttt{codes/run\_ablation.py}), this patch happens
automatically inside \texttt{main\_mmhcl\_plus.py}; no manual cell is required.

\subsection{Execution convention}
A single variant can be launched from the command line using:
\begin{lstlisting}[language=bash]
ABLATION=A4_no_ramp python train.py --use_mmhcl_plus True
\end{lstlisting}

\section{Ablation aggregation and marginal contribution}
The notes also describe exporting a CSV summary of the ablation results and computing the marginal contribution of each variant relative to \texttt{A0\_full}. In practice, after loading the results into a dataframe \texttt{df}, the baseline row is indexed by \texttt{A0\_full}, and each metric column is converted into a delta column \(\Delta\) with respect to that baseline.

This makes the final table suitable for direct insertion into a Q1-style paper, because the table reports both absolute performance and the degradation or improvement induced by each architectural intervention.

\section{Statistical validation}
For Q1-level reporting, the notes explicitly require paired significance testing between \texttt{A0\_full} and every ablation variant over at least three random seeds.

\begin{lstlisting}[language=Python]
from scipy.stats import ttest_rel, wilcoxon
# df_seeds: rows = seeds, cols = variants, values = NDCG@20
for v in df_seeds.columns.drop("A0_full"):
    t, p_t = ttest_rel(df_seeds["A0_full"], df_seeds[v])
    w, p_w = wilcoxon(df_seeds["A0_full"], df_seeds[v])
    print(f"{v:15s} | Delta={df_seeds['A0_full'].mean()-df_seeds[v].mean():+.4f} "
          f"| t-test p={p_t:.4f} | Wilcoxon p={p_w:.4f}")
\end{lstlisting}

The rationale is straightforward: mean performance alone is insufficient for a high-quality publication, whereas paired statistical tests demonstrate whether each ablation effect is robust across seeds.

\section{Repository checklist}
A coherent implementation workflow for the repository is as follows:
\begin{enumerate}[leftmargin=1.6em]
\item Add \texttt{mmhcl\_plus/ablation/ablation\_config.py} as the canonical registry.
\item Add or replace \texttt{mmhcl\_plus/trainer/hybrid\_balancer.py} with the extended balancer.
\item Copy the six ablation cells into the notebook while preserving the working BPR, VICReg, and InfoNCE pipeline with seed 612550091.
\item Run variants from the CLI using the \texttt{ABLATION} environment variable.
\item Aggregate all metrics into a CSV table and compute marginal deltas against \texttt{A0\_full}.
\item Perform paired t-tests and Wilcoxon signed-rank tests over multiple seeds before reporting final conclusions.
\end{enumerate}

\section{Closing note}
This English LaTeX version consolidates the remaining material visible in \texttt{paste.txt} into a more coherent document structure. It preserves the main intent of the original notes: to provide a reproducible diagnostic and ablation framework for MMHCL-Plus that is suitable for implementation in a GitHub repository and for reporting in a Q1-style academic manuscript.

\end{document}
